\documentclass[aspectratio=169,11pt]{beamer}

\usetheme{Madrid}
\usecolortheme{default}
\usefonttheme{professionalfonts}
\setbeamertemplate{navigation symbols}{}
\setbeamertemplate{footline}[frame number]

\usepackage{amsmath, amssymb, booktabs, graphicx}

\title{Culture, Norms, and Labor-Market Allocation}
\subtitle{Labor Markets and Political/Cultural Institutions --- Week 5}
\author{Teaching deck draft}
\date{}

\begin{document}

\begin{frame}
  \titlepage
\end{frame}

\begin{frame}{Central question}
\begin{itemize}
    \item How do norms about work, family, reputation, competition, leadership, and acceptable conduct shape labor allocation?
    \item Which margins move: labor supply, search, care work, entry, entrepreneurship, migration, promotion, or wage acceptance?
    \item How can economists identify norms rather than using culture as a residual label?
\end{itemize}
\end{frame}

\begin{frame}{Bridge from Week 4}
\begin{itemize}
    \item Week 4: worker voice institutions aggregate claims and change bargaining, governance, and spillovers.
    \item Week 5: norms shape the acceptable action set before formal organization may exist.
    \item Norms can alter who searches, which jobs are legitimate, who enters competition, and what wages can be accepted.
    \item Keep the same discipline: labor outcome, institutional object, mechanism, evidence.
\end{itemize}
\end{frame}

\begin{frame}{Norms as labor-market mechanisms}
\begin{itemize}
    \item A norm is a shared expectation with social, reputational, or identity consequences.
    \item It is not a residual after wages, law, and technology fail to explain an outcome.
    \item Norms enter decisions through stigma, approval, family obligation, status, religious reputation, and peer discipline.
    \item The empirical task is to separate norms from prices, legal constraints, discrimination, selection, and local demand.
\end{itemize}
\end{frame}

\begin{frame}{A compact allocation map}
\[
U_{ij} = w_{ij} - c_{ij} - s_{ij}(a_{ij}, n_i) + \varepsilon_{ij}
\]
\begin{itemize}
    \item $w_{ij}$: wage or material return.
    \item $c_{ij}$: direct cost of work, search, commuting, care, or mobility.
    \item $s_{ij}(a_{ij}, n_i)$: social cost or reward from action $a_{ij}$ under norm environment $n_i$.
    \item Norms change allocation when the social term shifts the net return to a labor-market action.
\end{itemize}
\end{frame}

\begin{frame}{Week 5 versus Week 6}
\small
\begin{tabular}{p{0.32\linewidth} p{0.29\linewidth} p{0.27\linewidth}}
\toprule
Object & Week 5 & Week 6 \\
\midrule
Main mechanism & conduct and role expectations & structural hierarchy and category-based access \\
Examples & family roles, reputation, competition, networking, work ethic & caste, race, ethnicity, migrant status, institutionalized group hierarchy \\
Labor margin & participation, search, entry, promotion entry, wage acceptance & hiring, authority, protection, networks, occupation boundaries \\
Diagnostic question & what conduct is rewarded or punished? & how does category membership structure access? \\
\bottomrule
\end{tabular}
\end{frame}

\begin{frame}{Family and gender-role norms}
\begin{itemize}
    \item Norms about mothers' work, breadwinner identity, care obligations, and acceptable occupations affect labor supply.
    \item Fernandez--Fogli: inherited culture and second-generation work and fertility outcomes.
    \item Alesina--Giuliano--Nunn: historical plough use and persistent gender-role norms.
    \item Jayachandran: norms as barriers to women's employment alongside prices, safety, childcare, and bargaining.
\end{itemize}
\end{frame}

\begin{frame}{Religion, reputation, and acceptable work}
\begin{itemize}
    \item Religion can operate through public visibility, mobility, sector choice, and reputation-preserving conduct.
    \item Carvalho: veiling as a conduct and reputation technology, not just a cultural label.
    \item The labor margins are participation, sector choice, entrepreneurship, mobility, and public-facing work.
    \item Boundary: religion as conduct norm belongs here; religion as durable group hierarchy belongs in Week 6.
\end{itemize}
\end{frame}

\begin{frame}{Competition and self-promotion}
\begin{itemize}
    \item Competitive framing can change who applies even when wage offers are held fixed.
    \item Flory--Leibbrandt--List: randomized job-entry features and application behavior.
    \item Delfino: gender-coded occupational entry and barriers to men entering pink-collar jobs.
    \item These designs identify entry responses to job framing, not the full social origin of norms.
\end{itemize}
\end{frame}

\begin{frame}{Leadership-entry and networking norms}
\begin{itemize}
    \item Advancement can depend on informal networking, visibility, manager interaction, and self-advocacy.
    \item Cullen--Perez-Truglia: schmoozing and the gender gap in career advancement.
    \item Week 5 question: who chooses or feels able to enter informal advancement systems?
    \item Week 6 question: who is structurally excluded from networks, authority, or evaluation?
\end{itemize}
\end{frame}

\begin{frame}{Work ethic and unemployment norms}
\begin{itemize}
    \item Work norms change the social utility of employment, unemployment, benefit use, and search.
    \item Clark: unemployment stigma varies with the local social environment.
    \item Eugster--Lalive--Steinhauer--Zweimuller: work attitudes and search near the Swiss language border.
    \item The observed margins are search, reservation behavior, benefit use, and psychological costs.
\end{itemize}
\end{frame}

\begin{frame}{Wage-acceptance norms}
\begin{itemize}
    \item Accepting a wage can have social consequences when workers monitor one another.
    \item Breza--Kaur--Krishnaswamy: social suppression of labor supply.
    \item The object is not only an individual reservation wage.
    \item The market-level mechanism is social discipline around undercutting and acceptable labor supply.
\end{itemize}
\end{frame}

\begin{frame}{Mechanism map}
\small
\begin{tabular}{p{0.30\linewidth} p{0.28\linewidth} p{0.28\linewidth}}
\toprule
Norm channel & Labor margin & Typical evidence \\
\midrule
Family roles & participation, hours, care & inherited culture, history \\
Reputation & public work, mobility, sector & religion and community models \\
Competition & applications, entry & job-entry experiments \\
Networking & promotion, leadership & workplace data \\
Work ethic & search, unemployment & local social environment \\
Wage acceptance & labor supply, reservation behavior & market-level experiments \\
\bottomrule
\end{tabular}
\end{frame}

\begin{frame}{Measurement strategies}
\small
\begin{tabular}{p{0.34\linewidth} p{0.28\linewidth} p{0.24\linewidth}}
\toprule
Design & Identifies best & Key limitation \\
\midrule
Inherited-culture design & transmission net of host environment & migrant selection, origin bundles \\
Historical persistence & origins and durability & correlated historical channels \\
Job-entry field experiment & application response to framing & partial norm environment \\
Workplace promotion data & informal advancement margins & norms versus access barriers \\
Local social environment & current social pressure & local demand and composition \\
Market-level labor experiment & social discipline in supply & setting-specific visibility \\
\bottomrule
\end{tabular}
\end{frame}

\begin{frame}{Diagnostics}
\begin{itemize}
    \item Is the proxy measuring norms rather than prices or constraints?
    \item Is the design identifying inherited transmission or current social pressure?
    \item Does the evidence observe conduct, beliefs, applications, promotions, wages, or search?
    \item What would the same pattern mean under discrimination, household bargaining, or structural hierarchy?
\end{itemize}
\end{frame}

\begin{frame}{Week 5 lab}
\begin{itemize}
    \item Reproduce: synthetic inherited-culture exercise inspired by Fernandez--Fogli.
    \item Diagnose: compare norm proxies with host prices, demand, childcare, education, and household controls.
    \item Transfer: classify job-entry, religion, networking, work-norm, and market-level designs.
    \item Core output: what proxy measures, what labor margin moves, and what the design does not identify.
\end{itemize}
\end{frame}

\begin{frame}{Frontier questions}
\begin{itemize}
    \item How do firm practices amplify or soften norms through job design, scheduling, remote work, and promotion rules?
    \item How fast do norms change when policy, technology, migration, or labor demand changes?
    \item When do community norms create local labor-supply distortions?
    \item How much of leadership gaps reflect conduct norms rather than structural hierarchy?
\end{itemize}
\end{frame}

\begin{frame}{Bridge to Week 6}
\begin{itemize}
    \item Week 5 studied expectations about conduct: work, family, reputation, competition, networking, and wage acceptance.
    \item Week 6 studies durable categories and institutionalized hierarchy.
    \item The transition question: when does a role norm become a system of access, exclusion, authority, and protection?
    \item Keep categories bounded in Week 5; use Week 6 for caste, race, ethnicity, migrant status, and structural hierarchy.
\end{itemize}
\end{frame}

\begin{frame}{Takeaways}
\begin{itemize}
    \item Norms change labor allocation through social costs, rewards, acceptable job sets, and community discipline.
    \item The key margins are participation, search, entry, promotion, migration, entrepreneurship, and wage acceptance.
    \item Credible evidence names whether it identifies transmission, current pressure, job-entry response, workplace advancement, or market-level supply discipline.
    \item Week 5 is conduct and role norms; Week 6 is structural hierarchy.
\end{itemize}
\end{frame}

\end{document}
